% !PRINTED 690
\weberpage{690}

Every differential $dJ$ whose lower polygon is $\mathfrak{B}$, or a divisor
of $\mathfrak{B}$, can then be represented in the form:
\weq{7}{dJ = c_1dJ_1 + c_2dJ_2 + c_3dJ_3 + \cdots}
with constant coefficients.

The simplest differentials are those whose lower polygon is the
null polygon, and whose upper polygons $\mathfrak{W}$ are accordingly of order
$(2p - 2)$. These are called \gesperrt{differentials of the first kind}
and are denoted by $dW$. They are characterized by the fact
that for every variable $z$ the polygons $\mathfrak{U}^2\mathfrak{W}$ and $\mathfrak{Z}$ must be
equivalent, and the $\mathfrak{W}$ accordingly form, assuming their existence, a
class of polygons $W$ whose dimension is to be determined.

The polygon $\mathfrak{W}$ is called the fundamental polygon of the differential $dW$
and is named a \gesperrt{complete polygon of the first kind}. If
$\mathfrak{W} = \mathfrak{A}\mathfrak{B}$, then $\mathfrak{A}$ and $\mathfrak{B}$ too are called polygons of the first
kind, and indeed \gesperrt{complementary polygons} of one another.

Every polygon that is not a divisor of a complete polygon
of the first kind is said to be \gesperrt{of the second kind}.

If $\mathfrak{A}$ is a polygon of the first kind, then every polygon $\mathfrak{A}'$
equivalent to $\mathfrak{A}$ is also of the first kind, and we call
the class $A$ of $\mathfrak{A}$ a \gesperrt{class of the first kind}. For if
$\mathfrak{A}\mathfrak{B} = \mathfrak{W}$, then $AB = W$, where $B$ is the class of $\mathfrak{B}$, and
$\mathfrak{A}'\mathfrak{B} = \mathfrak{W}'$ is equivalent to $\mathfrak{W}$, hence of the first kind.

If, then, $q$ is the number of linearly independent polygons
of the first kind that are divisible by a polygon of the class $A$,
then by § 192, (14)
\weq{8}{q = (A,\ W) = (O,\ B),}
that is, equal to the dimension of the complementary class.

If $A$ is a class of the second kind, then there is no polygon
of the first kind divisible by $\mathfrak{A}$, and
\[ (A,\ W) = 0. \]

\begin{center}
§ 198. \textbf{\gesperrt{The Family of Polygons of the First Kind.}}
\end{center}

In order to determine the dimension $(O, W)$ of the family of polygons of the first kind and
thereby at the same time to establish its existence, we take
an arbitrary variable $z$ in $\Omega$ with the branch polygon
$\mathfrak{Z}$ and the lower polygon $\mathfrak{U}$. If $dW$ is a differential of the first kind,
then
\weq{1}{w = \frac{dW}{dz} = \frac{\mathfrak{W}\mathfrak{U}^2}{\mathfrak{Z}}}
%CONT>>
